\begin{textsample}{POS Dim 2 – am – Score 23.00 – am/am\_em\_28.txt}  \label{ex:f2_pos_155}
3.3.2 Unidades Curriculares de Aprofundamentos ( UCAs ) As Unidades Curriculares de Aprofundamento têm por objetivo aprofundar as aprendizagens relacionadas às competências gerais, às Áreas de Conhecimento e/ou à Formação \textbf{Técnica} e Profissional.
No caso dos \textbf{Itinerários} \textbf{Formativos} de Áreas do Conhecimento, essa ampliação acontece em articulação com temáticas contemporâneas \textbf{sintonizadas} com o contexto e os interesses dos estudantes.
Já na Formação \textbf{Técnica} e Profissional, a expansão se dá juntamente com o desenvolvimento de habilidades básicas requeridas pelo mundo do trabalho e habilidades específicas relacionadas aos \textbf{Cursos} \textbf{Técnicos}, \textbf{Cursos} de \textbf{Qualificação} \textbf{Profissional} ( FICs ) ou Programa de Aprendizagem escolhidos pelos estudantes.
Por Programa de Aprendizagem, segundo Resolução CNE/CEB nº 3/2018, compreende -se “ arranjos e combinações de \textbf{cursos} que, articulados e com os devidos aproveitamentos curriculares, possibilitam um \textbf{itinerário} \textbf{formativo} ” ( BRASIL, 2018⁴ ).
Recomenda -se que as UCAs relacionadas às Áreas do Conhecimento incorporem e integrem os quatro eixos \textbf{estruturantes}, mencionados anteriormente.
Nas UCAs relacionadas à Formação \textbf{Técnica} e Profissional, essa incorporação fica a critério da \textbf{rede} e instituição \textbf{ofertante}, ficando a obrigatoriedade voltada aos eixos tecnológicos definidos pelo \textbf{Catálogo} Nacional de \textbf{Cursos} \textbf{Técnicos} ( CNCT ).
Além disso, as UCAs podem ser organizadas por diferentes arranjos curriculares : por disciplinas, por \textbf{oficinas}, por unidades/campos temáticos, por projetos, entre outras possibilidades de flexibilização.
Devem ser estruturadas e organizadas em etapas com terminalidade ( semestral ), possibilitando, assim, a emissão de \textbf{certificação} de \textbf{qualificação} para o trabalho, como prevê a Resolução CNE/CEB nº 3/2018.
As redes/instituições terão autonomia para definir a melhor estratégia pedagógica para \textbf{oferta} das UCAs.
Entretanto, devem obedecer à distribuição de \textbf{carga} \textbf{horária} e garantir a \textbf{oferta} semestral prevista no Quadro 17, de modo a não comprometer a vida escolar dos estudantes em caso de mudança de \textbf{rede} e/ou instituição de ensino.
Vale destacar que a construção dos Aprofundamentos deve ter por base um bom diagnóstico que considere : os interesses dos estudantes, por meio da escuta, as potenciais parcerias e a condição de \textbf{oferta} das escolas e da própria \textbf{rede} de ensino. 3.3.3 Unidades Curriculares \textbf{Eletivas} ( UCEs ) Por Unidades Curriculares \textbf{Eletivas} entendem -se as estratégias curriculares que possibilitam a experimentação por parte do estudante dos conhecimentos das diversas áreas do conhecimento, podendo estar atreladas ou não à área de conhecimento escolhida para aprofundamento.
Entretanto, não devem ser confundidas com componentes tradicionais nem como apoio pedagógico ou reforço escolar.
Por serem unidades curriculares de livre escolha, devem possibilitar aos estudantes a experimentação de diferentes temas, vivências e aprendizagens, de maneira a diversificar e enriquecer seu Itinerário \textbf{Formativo}.
Apesar do seu caráter mais lúdico e prático, é importante que as \textbf{eletivas} tenham intencionalidade pedagógica e se articulem com as Áreas do Conhecimento, os eixos \textbf{estruturantes} e as Competências Gerais da BNCC.
Na Formação \textbf{Técnica} e Profissional, a \textbf{oferta} da UCE deve ocorrer por meio de uma FIC.
Assim como ocorre com a \textbf{oferta} das demais unidades curriculares que compõem o IF, as redes/instituições terão autonomia para definir a melhor estratégia pedagógica para \textbf{oferta} das UCEs.
Entretanto, devem obedecer à distribuição de \textbf{carga} \textbf{horária} e garantir a \textbf{oferta} semestral conforme Quadro 17, evitando prejuízos à vida escolar dos estudantes em caso de mudança de \textbf{rede} e/ou instituição de ensino.
Na 2º série do ensino médio deve ser garantida a \textbf{oferta} de uma \textbf{eletiva} de 40 horas por \textbf{semestre}.
Na 3º série deve ser garantida a \textbf{oferta} de 80 horas de \textbf{eletiva} por \textbf{semestre}, podendo ser dividida em duas unidades curriculares de 40 horas por \textbf{semestre}, totalizando quatro \textbf{eletivas} ao final do ano escolar, ou uma \textbf{eletiva} de 80 horas por \textbf{semestre}, totalizando duas UCEs.
Como se pode perceber, devem ser \textbf{ofertadas} pelo menos duas \textbf{eletivas} ao longo de um ano, com \textbf{carga} \textbf{horária} que pode variar em uma ou duas aulas \textbf{semanais}.
Ressalta -se que, quando a \textbf{oferta} da \textbf{Eletiva} ocorrer por meio de uma FIC, cuja \textbf{carga} \textbf{horária} mínima é de 160 horas, a \textbf{oferta} poderá se estender por mais de um \textbf{semestre}.
Recomenda -se que as Unidades Curriculares \textbf{Eletivas} \textbf{atendam} aos \textbf{requisitos} mínimos exigidos aos componentes de eletividade, de acordo com a ementa a seguir.
Quadro 18 - Sugestão para a construção de Ementas de Unidade Curricular \textbf{Eletiva} - Título : nome objetivo e atraente que facilite a compreensão e motive a escolha dos estudantes. - Proponente : nome da escola que propõe a UCE. - Professor(es) Responsável(eis) : nome do(s) professor(es) autores da \textbf{Eletiva}. - Resumo : descrição sucinta e interessante que ajude professores e estudantes a compreenderem a proposta da UCE. - Área de Conhecimento : indicação da(s) Área(s) do Conhecimento a serem trabalhadas pela \textbf{Eletiva}, lembrando a recomendação de que sejam interdisciplinares e possam aprofundar e ampliar aprendizagens em uma ou mais Áreas do Conhecimento.
No caso específico do Itinerário de Formação \textbf{Técnica} e \textbf{Profissional}, considerar os eixos do \textbf{Catálogo} Nacional de \textbf{Cursos} \textbf{Técnicos} ( CNCT ) e a Classificação Brasileira de Ocupações ( CBO ). - Habilidades : indicação das habilidades a serem desenvolvidas, lembrando que as \textbf{Eletivas} podem ter diversos formatos e abordar diferentes objetos de conhecimento desde que trabalhem de forma intencional as aprendizagens relacionadas às Áreas do Conhecimento, às Competências Gerais da BNCC ou a, pelo menos, um eixo \textbf{estruturante} dos \textbf{Itinerários} \textbf{Formativos}.
No caso específico do Itinerário de Formação \textbf{Técnica} e Profissional, é preciso também considerar os eixos do \textbf{Catálogo} Nacional de \textbf{Cursos} \textbf{Técnicos} ( CNCT ) e a Classificação Brasileira de Ocupação ( CBO ). - Objetos de Conhecimento : identificação dos objetos de conhecimento a serem estudados ao longo da \textbf{Eletiva}. - Eixos \textbf{Estruturantes} : indicação de que ( quais ) eixo(s) será(ão) trabalhado(s) pela \textbf{Eletiva}. - Objetivos : descrição das mudanças que se espera promover nos estudantes. - Unidade Curricular : definição da natureza da \textbf{Eletiva} ( núcleo de estudos, laboratório, projeto, \textbf{oficina}, FIC, dentre outros ). - Sequência de Situação/Atividade Educativas : roteiro de estratégias metodológicas. - \textbf{Carga} \textbf{Horária} : indicação da duração de cada eixo \textbf{estruturante} e/ou de cada situação ou atividade educativa. - \textbf{Perfil} Docente : indicação de quantos professores serão necessários, bem como dos conhecimentos, habilidades e características que eles devem ter. - \textbf{Perfil} dos Participantes : indicação de faixa etária, ano, interesses, além de quantidade mínima e máxima de estudantes por turma. - Recursos : indicação dos espaços, equipamentos e materiais necessários. - Avaliação : definição de como avaliar o desenvolvimento dos estudantes. - Fontes de informação : indicação de livros, sites, vídeos e outros materiais de referência para subsidiar o trabalho com a \textbf{Eletiva}.
Fonte : Adaptado de Brasil, 2019.
A sugestão descrita no Quadro 18 contempla os \textbf{requisitos} mínimos para a \textbf{oferta} de uma UCE, de maneira que compete a cada \textbf{rede} a ampliação do ementário, de acordo com suas propostas pedagógicas.

% matched lemmas: atender, carga, catálogo, certificação, curso, eletivo, estruturante, formativo, horário, itinerário, oferta, ofertante, ofertar, oficina, perfil, profissional, qualificação, rede, requisito, semanal, semestre, sintonizar, técnico
\end{textsample}
